\section{Time domain: attack, overshoot, settling; the driver invariants}
\label{sec:time}

Model a bass attack as a one-cycle acceleration burst, $\omega\gg\wc$
(free-mass limit).

\begin{lemma}[Burst shape factor]
\label{lem:shape}
With $C\equiv \omega^2\max_t|x(t)|/A$: $C=1$ for steady sine; for one-cycle
sine-phase bursts $C=2\pi$, cosine-phase $C=2$; $n$ cycles, sine phase,
$C=2\pi n$. Transient program thus demands $C\in[2,2\pi]$
($6$--$16$\,dB) more excursion per unit peak pressure than steady sine.
\end{lemma}

Peak-equivalent ceilings become
\begin{align}
\hat L_{\mathrm{L,burst}}(f)&=108.5+20\log_{10}(f^2\Vd)-20\log_{10}C,
\label{eq:burstL}\\
\hat L_{\mathrm{T,burst}}&=112.2+10\log_{10}(\eta_0\,\kappa\Pmax),
\label{eq:burstT}
\end{align}
with $\kappa$ (A5) the short-term dissipation headroom.

\begin{proposition}[Equivalence of the two analyses; boundary shift]
\label{prop:equiv}
Time- and frequency-domain limits yield the same two figures of merit,
$\Vd$ and $\eta_0\Pmax$, rescaled by $C$ and $\kappa$. The regime boundary
moves up to
\begin{equation}
\hat f_x=f_x\,(\kappa C^2)^{1/4}
\qquad(\kappa=4,\ C=2\ \Rightarrow\ \hat f_x=2f_x).
\label{eq:fxburst}
\end{equation}
\end{proposition}

So for attack-type program the excursion regime extends roughly an octave
higher than the sine analysis suggests; both roles need both FOMs checked at
their own band edge.

\begin{proposition}[Settling is a driver constant; ring shape is a box choice]
\label{prop:decay}
The envelope decay rate of every sealed alignment of a given driver equals
$\sigma/2=\pi F_s/\Qts$ (box-invariant, \cref{prop:invariance}); the box
selects only $\omega_d$ and the overshoot pattern. Step-response overshoot
$M_p=\exp(-\pi\zeta/\sqrt{1-\zeta^2})$, $\zeta=1/(2\Qtc)$: $M_p\le1\%
\Leftrightarrow\Qtc\le0.61$; $M_p\le5\%\Leftrightarrow\Qtc\le0.72$.
\end{proposition}

Spurious ring-down tail pressure satisfies
$\hat p_{\mathrm{tail}}/\hat p_{\mathrm{burst}}\le C(\wc/\omega)^{2}$
(worst phase), improving 12\,dB/oct with corner-below-band separation --- an
LTI, hence FIR-correctable, effect (at the cost of \eqref{eq:EQtax}). The
invariants FIR cannot touch remain $\Vd$, $\eta_0\Pmax$, $\sigma$.

\subsection{The driver invariants, and what \texorpdfstring{$EBP$}{EBP} actually is}
\label{sec:invariants}

Four quantities survive both analyses unchanged: $\Vd=\Sd\Xmax$ (excursion
regime); $\eta_0\Pmax$ (dissipation regime, \eqref{eq:eta0}); the corner
rate $\sigma/2\pi=F_s/\Qts$ (placement, settling); and $f_L=\Re/(2\pi\Le)$
(upper validity). Cone area $\Sd$ enters again separately, via Doppler
(\cref{sec:distortion}).

A common rule of thumb states $EBP\propto V_c/\Mms$ as a motor-strength
proxy; here is the exact version. For an overhung motor with coil conductor
mass $m_c$, conductivity $\sigma_c$, density $\rho_c$, flux density $B$, and
effective overhang factor $u\ge1$:
\begin{equation}
\se=\frac{\Bl^2}{\Re\Mms}=\frac{B^2\sigma_c}{\rho_c}\cdot\frac{m_c}{\Mms}\cdot\frac{1}{u^2}.
\label{eq:sebound}
\end{equation}

\begin{theorem}[Motor materials bound]
\label{thm:bound}
For an overhung motor, equality holds in \eqref{eq:sebound}:
$EBP\cdot u^{2}=K_{\mathrm{mat}}\cdot\dfrac{m_c}{\Mms}$, with
$K_{\mathrm{mat}}\equiv\dfrac{B^{2}\sigma_c}{2\pi\rho_c}$ (copper:
$1059\,B^2$\,Hz; aluminium: $2063\,B^2$\,Hz). Bounding the coil-mass
fraction, $m_c/\Mms\le\beta_{\max}\lesssim0.35$, gives
\begin{equation}
\boxed{\ EBP\cdot u^{2}\;\le\;K_{\mathrm{mat}}\,\beta_{\max}\ }
\label{eq:thmbound}
\end{equation}
\end{theorem}

At $B=1.0$--$1.2$\,T (Cu, $\beta_{\max}=0.35$) the bound is
$\approx370$--$530$\,Hz: a sub-class driver ($u\approx3$) is limited to
$EBP\lesssim45$\,Hz, an attack-class corner rate of $135$\,Hz forces
$u\lesssim1.8$. A single driver combining sub-class stroke and attack-class
corner rate would need $\beta B^2\approx1.1$ --- impossible. Three
independent mechanisms enforce the split: \eqref{eq:thmbound} itself
(stroke suppresses corner rate quadratically); \cref{prop:invariance} (one
driver has one $\sigma$, hence one $F_c$); and A2 (tall coil $\Rightarrow$
large $\Le$ $\Rightarrow$ low $f_L$ intruding into the attack band).
Escape routes (underhung/XBL$^2$, dual gap, shorting rings, aluminium
coils) buy factors of $\sim2$, not the required $\sim5$, trading against
$\beta$ or $B$ inside the same bound.

$\beta_{\max}$ is an observed ceiling on real motor construction (coil
formers, adhesives, winding density), not a law of electromagnetism, and
neither $m_c$ nor $B$ is a standalone datasheet field. Since $u\ge1$,
however, \eqref{eq:thmbound}'s equality also gives a lower bound
computable from $EBP$ alone: $m_c/\Mms\ge EBP/K_{\mathrm{mat}}(B)$,
independent of the otherwise-unmeasurable $u$. \Cref{tab:ebp-census}
applies this to $N=1025$ bass-relevant drivers (VituixCAD database, Type
S/W/WM): at a conservative $B=1.0$\,T (copper), only $0.3\%$ approach
$\beta_{\max}$, and the material-agnostic (aluminium) bound is never
exceeded; the single outlier needs only $B\approx1.4$\,T (a plausible
neodymium motor) to reach the ceiling, not exceed it.

\begin{table*}[htbp]
\centering
\caption{Lower-bound coil-mass fraction $\beta_{\min}=EBP/K_{\mathrm{mat}}(B)$ (Thm.~\ref{thm:bound} at $u{=}1$, Cu $B{=}1.0$\,T unless noted) over $N{=}1025$ bass-relevant drivers (Type S/W/WM) in the bundled VituixCAD database, vs.\ $\beta_{\max}\approx0.35$.}
\label{tab:ebp-census}
\small
\begin{tabular}{lc}
\toprule
quantity & value\\
\midrule
$N$ & 1025\\
median $\beta_{\min}$ & 0.10\\
$p_{90}/p_{99}/p_{99.5}$ & 0.17 / 0.29 / 0.31\\
max & 0.66 (B\&C Speakers 8NSM64, needs $B{\approx}1.37$\,T)\\
material-agnostic (Al) excess & 0.0\%\\
\bottomrule
\end{tabular}
\end{table*}

